% Part: lambda-calculus
% Chapter: syntax
% Section: eta

\documentclass[../../../include/open-logic-section]{subfiles}

\begin{document}

\olfileid{lam}{syn}{eta}
\olsection{$\eta$-转换}

在 $\lambda$ 项上还存在另一种关系。在 \olref[fv]{sec} 节中，
我们使用了 $\lambd[x][(fx)]$ 这个例子，它接受一个自变元并对之应用 $f$。
换句话说，它与~$f$ 是同一个函数：$\lambd[x][(fx)]N$ 和 $fN$ 都归约到~$fN$。
我们用 $\eta$-归约（以及 $\eta$-扩张）来刻画这一思想。

\begin{defn}[$\eta$-收缩, $\eredone$]
  \ollabel{defn:beredone}
  \emph{$\eta$-收缩} ($\eredone$) 是指项上满足以下条件的最小兼容关系：
  \[
    \lambd[x][M x] \eredone M \text{条件是} x \notin FV(M)
  \]
\end{defn}

\begin{defn}[$\beta\eta$-归约, $\bered$] \ollabel{defn:bered}
  \emph{$\beta\eta$-归约} ($\bered$) 是指项上包含 $\bredone$ 和 $\eredone$ 的最小自反且传递的关系，
  即自反性、传递性规则加上以下两条规则：
  \begin{enumerate}
  \item 若 $M \bredone N$ 则 $M \bered N$。\ollabel{defn:bered3}
  \item 若 $M \eredone N$ 则 $M \bered N$。\ollabel{defn:bered4}
  \end{enumerate}
    
\end{defn}

\begin{defn}
  我们用 $\eta$-转换规则来扩展等价关系 $\equal$：
  \[
  \lambd[x][f x] \equal f
  \]
  并将扩展后的关系记为 $\equal[\eta]$。
\end{defn}

$\eta$等价之所以重要，是因为它与 $\lambda$ 项的外延性有关：

\begin{defn}[外延性]
  我们用 (\ext) 规则来扩展等价关系 $\equal$：
  \begin{center}
  若 $Mx \equal Nx$ 则 $M \equal N$，条件是 $x \notin FV(MN)$。
  \end{center}
  并将扩展后的关系记为 $\equal[\ext]$。
\end{defn}

粗略地说，这条规则表明，若将两个项视为函数，且它们对相同的自变元表现相同，则应被视为相等。

现在我们来证明，$\eta$ 规则恰好提供了外延性，且不提供其他内容。

\begin{thm}
  $M \equal[\ext] N$ 当且仅当 $M \equal[\eta] N$。
\end{thm}

\begin{proof}
  首先我们证明 $\equal[\eta]$ 在外延性规则下是封闭的。
  也就是说，\ext 规则不会对 $\equal[\eta]$ 增加任何新内容。
  这样一来，我们就知道 $\equal[\eta]$ 包含 $\equal[\ext]$，
  并且如果 $M \equal[\ext] N$，那么 $M \equal[\eta] N$。

  为证 $\equal[\eta]$ 在 \ext 规则下封闭，注意到对于由 \ext 规则推导出的任何 $M \equal N$，
  其前提是 $Mx \equal[\eta] Nx$。于是根据 $\equal$ 的规则有 $\lambd[x][Mx]
  \equal[\eta] \lambd[x][Nx]$，对两边应用 $\eta$
  便得到 $M \equal[\eta] N$。

  类似地，我们证明 $\eta$ 规则包含在 $\equal[\ext]$ 中。
  对于任何满足 $x \notin FV(M)$ 的 $\lambd[x][Mx]$ 和 $M$，我们有 $(\lambd[x][Mx])x \equal[\ext] Mx$，
  于是通过 \ext 规则得到 $\lambd[x][Mx] \equal[\ext] M$。
\end{proof}

\end{document}